% =============================================================================
% IEEE TVLSI Paper: A 12-bit 10-MS/s SAR ADC With Current-Domain
%                     kT/C Noise Cancellation (CAAZ)
% Compile: pdflatex -> bibtex -> pdflatex -> pdflatex
% =============================================================================
\documentclass[journal]{IEEEtran}

\usepackage{amsmath,amssymb}
\usepackage{algorithmic}
\usepackage{array}
\usepackage[caption=false,font=normalsize,labelfont=sf,textfont=sf]{subfig}
\usepackage{textcomp}
\usepackage{url}
\usepackage{verbatim}
\usepackage{graphicx}
\usepackage{cite}
\usepackage{booktabs}
\hyphenation{ana-log-to-dig-i-tal semi-conduc-tor}
\def\BibTeX{{\rm B\kern-.05em{\sc i\kern-.025em b}\kern-.08em
    T\kern-.1667em\lower.7ex\hbox{E}\kern-.125emX}}
\usepackage{balance}
\usepackage{siunitx}

\graphicspath{{../05_研究报告/figures/}}

\begin{document}

\title{A 12-bit 10-MS/s SAR ADC With Current-Domain $kT/C$ Noise Cancellation via CAAZ}

\author{Zhao Yi,~\IEEEmembership{Student Member, IEEE}
\thanks{This work was supported by the National Natural Science Foundation of China under Grant XXXXXXX.}
\thanks{Z. Yi is with the School of Integrated Circuits, Huazhong University of Science and Technology, Wuhan 430074, China (e-mail: yizhao@hust.edu.cn).}}

\markboth{IEEE Transactions on Very Large Scale Integration (VLSI) Systems}%
{Yi\MakeLowercase{\textit{et al.}}: A 12-bit 10-MS/s SAR ADC With Current-Domain $kT/C$ Noise Cancellation via CAAZ}

\IEEEpubid{XXXX-XXXX/XX\$XX.XX~\copyright~2026 IEEE}

\maketitle

\begin{abstract}
This paper presents a 12-bit 10-MS/s SAR ADC employing a current-adjusting auto-zeroing (CAAZ) technique for $kT/C$ noise cancellation in 180-nm CMOS. The CAAZ architecture overcomes the fundamental gain-saturation trade-off of conventional voltage-domain output offset storage (OOS) approaches by shifting noise extraction to the current domain. A capacitance burden shifting mechanism relocates the storage capacitor to a low-impedance internal node, enabling $C_{\rm AZ}=300$~fF to suppress secondary noise while keeping input capacitance minimal. The CAAZ noise cancellation efficiency $\eta$ is experimentally characterized using a 100$\times$ noise injection (NF=10) double-$t_d$ method across five PVT corners. Transient noise simulations demonstrate an ENOB of 11.61~bit and SNR of 71.81~dB at the TT corner (NF=1, $t_d$=26~ns), with all corners exceeding 11.4~bit. The CAAZ $\eta$ reaches 87.9\% at TT, reducing to 79.7\% at the worst-case SF corner (85$^\circ$C, slow-N fast-P). Total power consumption is 720~$\mu$W from a 1.8-V supply. A two-batch cross-validation confirms $\eta$ consistency within 1.4~p.p.
\end{abstract}

\begin{IEEEkeywords}
SAR ADC, $kT/C$ noise cancellation, current-adjusted auto-zeroing, CAAZ, 180-nm CMOS.
\end{IEEEkeywords}

\section{Introduction}
\IEEEPARstart{S}{uccessive} approximation register (SAR) analog-to-digital converters (ADCs) dominate medium-resolution medium-speed applications owing to their excellent power efficiency~\cite{liu12Bit10MS2016}. Beyond 12-bit resolution, the fundamental $kT/C$ sampling noise---$v_{n,kT/C}=\sqrt{kT/C}$---emerges as the primary SNR bottleneck~\cite{kapustaSamplingCircuitsThat2014}. Suppressing this noise in a 12-bit system with 1.8-V reference to below 128~$\mu$V$_{\rm rms}$ demands picofarad-range sampling capacitors, incurring severe power and area penalties on the input driver. Conventional designs have used capacitors as large as 15~pF~\cite{liu12Bit10MS2016}, resulting in prohibitive front-end power.

Liu~\textit{et al.}~\cite{liu13bit0005mm240MS2020} pioneered voltage-domain $kT/C$ noise cancellation for SAR ADCs, reducing the sampling capacitor to 240~fF. However, this approach exposes a fundamental gain-saturation trade-off inherent to output offset storage (OOS): the preamplifier gain $A$ both suppresses post-amplifier noise and amplifies the differential input toward the supply rail. For a 12-bit converter with 0.9-V MSB step, even $A=10$ produces 9-V output, forcing $A\approx6$ and capping $kT/C$ attenuation at $\sim$16~dB~\cite{liu13bit0005mm240MS2020}. Subsequent efforts---PSD-bandwidth decoupling~\cite{li2020_SAR}, feedback auxiliary amplifiers~\cite{won2024_13b}, pre-sampling~\cite{guImprovedKTNoise2024}, and split sampling~\cite{huang5MS16bitLownoise2025}---each address specific aspects but introduce alternative burdens: complex timing, larger area, or removal of the preamplifier from the signal path. The common root cause is the OOS storage capacitor lying within the signal path, coupling noise suppression to distortion.

Safiallah~\textit{et al.}~\cite{safiallah2023_Currentadjusting} proposed a current-adjusting auto-zeroing (CAAZ) scheme that shifts offset storage outside the signal path, performing noise extraction in the current domain at a low-impedance internal node. This work builds upon that principle to deliver a complete 12-bit 10-MS/s SAR ADC in 180-nm CMOS. The principal contributions are:
\begin{enumerate}
    \item A current-domain CAAZ architecture achieving $\eta=87.9\%$ noise cancellation efficiency at TT corner, characterized via a 100$\times$ noise injection (NF=10) double-$t_d$ methodology across five PVT corners---the first such comprehensive characterization for CAAZ-based SAR ADCs.
    \item A capacitance burden shifting mechanism enabling $C_{\rm AZ}=300$~fF at a low-impedance internal node, reducing input capacitance while fully suppressing secondary noise.
    \item Transient noise simulations achieving ENOB=11.61~bit and SNR=71.81~dB at the TT corner, with all five PVT corners exceeding 11.4~bit, and total power of 720~$\mu$W. Two-batch cross-validation confirms $\eta$ consistency within 1.4~p.p.
\end{enumerate}

\section{Background and CAAZ Architecture}
\label{sec:caaz}

\subsection{Limitations of Voltage-Domain OOS}

The conventional voltage-domain $kT/C$ noise cancellation~\cite{liu13bit0005mm240MS2020} stores the amplified noise on a feedback capacitor during sampling and subtracts it during conversion. The residual input-referred noise is:
\begin{equation}
v_{\rm ns,res}=v_{\rm ns1}e^{-2\Delta t/\tau}+\frac{v_{\rm ns,AZ}}{A}+\frac{v_{\rm ns,comp}}{A}
\label{eq:oos}
\end{equation}
where $v_{\rm ns1}$ is the main sampling noise, $\Delta t$ the noise-sampling window, $\tau$ the settling constant, and $A$ the preamplifier gain. The exponential term demands $\Delta t\gg\tau$. The $1/A$ terms reveal that secondary noise is suppressed only by $A$, which is bounded by output swing.

The output swing limitation is decisive. With $V_{\rm REF}=1.8$~V, the differential MSB step reaches $\Delta V_{\rm in}\approx0.9$~V. Gain $A=10$ produces 9-V output, far beyond the 1.8-V supply, compelling $A\approx6$ that attenuates secondary noise by only $1/36$ (15.6~dB). Total $kT/C$ suppression is confined to $\sim$16~dB~\cite{liu13bit0005mm240MS2020}.

\subsection{CAAZ Operating Principle}

The CAAZ technique~\cite{safiallah2023_Currentadjusting} resolves this dilemma by transferring noise extraction to the current domain. Fig.~\ref{fig:caaz} depicts the proposed preamplifier with three-phase operation.

\begin{figure}[htbp]
\centering
\includegraphics[width=0.95\columnwidth]{Fig6_Preamplifier.png}
\caption{Proposed preamplifier with CAAZ technique.}
\label{fig:caaz}
\end{figure}

\textit{Phase 1 -- Main Sampling:} The input switch opens, freezing $V_{\rm in}$ with $kT/C$ noise $v_{\rm ns1}$ on $C_1$.

\textit{Phase 2 -- Current-Domain Noise Sampling:} Preamplifier inputs short to $V_{\rm CM}$ ( $\Delta V_{\rm in}=0$). Loads $M_{3,4}$ are diode-connected, limiting gain to $A_{\rm samp}\approx1$ to prevent saturation. Differential voltage $V_{\rm diff}=-v_{\rm ns1}$ converts to current stored on $C_{\rm AZ}$:
\begin{equation}
I_{\rm diff}=g_{mp}V_{\rm diff}=-g_{mp}v_{\rm ns1}
\label{eq:I_diff}
\end{equation}

\textit{Phase 3 -- SAR Conversion:} Loads switch to current-source mode, releasing $A_{\rm amp}=g_{mp}R_{\rm OA}>50$. The stored current cancels $v_{\rm ns1}$ while high gain suppresses post-amplifier noise.

\subsection{Current-Domain Noise Derivation}

Using correlated double sampling (CDS) theory, Table~\ref{tab:node} summarizes key node voltages.

\begin{table}[htbp]
\centering
\caption{Key Node Voltages in CAAZ Operation}
\label{tab:node}
\small
\begin{tabular}{lp{4.8cm}}
\toprule
\textbf{Phase} & \textbf{Voltage Expression} \\
\midrule
$t_1$ (Sampling) & $V_{C1}=V_{\rm in}(t_1)+v_{\rm ns1}$ \\[3pt]
$t_2$ (Amplify)  & $V_{C2}=g_{mp}R_{\rm out}[V_{\rm in}(t_2)-V_{\rm in}(t_1)-v_{\rm ns1}-v_{\rm os}]+v_{\rm ns2}$ \\
\bottomrule
\end{tabular}
\par\footnotesize{$R_{\rm out}=R\parallel1/g_{mn}$; $v_{\rm ns2}$: noise at $t_2$; $v_{\rm os}$: input offset.}
\end{table}

The input pair converts differential voltage to $I_{\rm diff}=g_{mp}[V_{\rm in}(t_2)-V_{\rm in}(t_1)-v_{\rm ns1}-v_{\rm os}]$; the load injects $I_{\rm ns}=g_{mn}v_{\rm ns2}$. Converting to voltage and dividing by SAR-phase gain $g_{mp}R_{\rm OA}$ yields:
\begin{IEEEeqnarray}{rCl}
v_{\rm ns,eq}&\approx&\frac{v_{\rm ns1}}{g_{mp}R_{\rm OA}}
+\frac{v_{\rm ns2}}{g_{mp}R_{\rm OA}}\cdot\frac{g_{mn}}{g_{mp}}
\label{eq:noise_simple}
\end{IEEEeqnarray}

The first term is the \textbf{residual main noise} $v_{\rm ns1,res}=v_{\rm ns1}/(g_{mp}R_{\rm OA})$, scaling with intrinsic gain $g_{mp}R_{\rm OA}>50$. The second term is the \textbf{secondary noise} from $C_{\rm AZ}$, requiring $C_{\rm AZ}=300$~fF to directly attenuate $\sqrt{kT/C_{\rm AZ}}$.

\subsection{Capacitance Burden Shifting}

The defining CAAZ innovation is \textit{capacitance burden shifting}. In OOS, the storage capacitor sits at the high-impedance output ($R_{\rm out}\sim100$~k$\Omega$). In CAAZ, $C_{\rm AZ}$ resides at a low-impedance node ($1/g_{mn}\sim1$~k$\Omega$):
\begin{equation}
\tau_{\rm external}=R_{\rm driver}C_1\;\gg\;\tau_{\rm internal}=C_{\rm AZ}/g_{mn}
\label{eq:tau}
\end{equation}
Thus $C_1$ is minimized, while $C_{\rm AZ}=300$~fF charges within $\sim$1~ns. The auxiliary amplifier operates at $\sim$1\% duty cycle, making its power overhead negligible.

\section{Circuit Implementation}
\label{sec:imp}

The preamplifier uses a pMOS input pair with nMOS active loads. During Phase~2, diode-connected loads provide $\approx$1 gain, preventing saturation under 0.9-V MSB steps. During SAR conversion, current-source mode releases $g_{mp}R_{\rm OA}>50$. $C_{\rm AZ}=300$~fF at the nMOS load gates stores current information as $V_{\rm GS}$ modulation. Dummy switches with complementary clocks suppress charge injection; bottom-plate sampling keeps errors below $1/4$~LSB. A StrongARM latch provides final comparison with kickback isolation. The optimal CAAZ settling time is $t_d=26$~ns, experimentally determined from the 24-26~ns sweep.

\section{Characterization Methodology}
\label{sec:method}

The CAAZ noise cancellation efficiency is characterized using a noise-factor (NF) injection method. The noise power injected into the sampling switches (I1 and I14) is amplified by $\text{NF}^2$ in the power domain. With NF=10, the injected noise power is 100$\times$ the thermal noise floor. The CAAZ cancellation efficiency $\eta$ is derived from the double-$t_d$ experiment:

\begin{equation}
\eta = 1 - \frac{101\cdot 10^{-\Delta\text{SNR}/10} - 1}{100}
\label{eq:eta}
\end{equation}

where $\Delta\text{SNR} = \text{SNR}_{t_d=26\text{ns}} - \text{SNR}_{t_d=24.05\text{ns}}$ (the shorter window approximates $\eta\approx0$). This methodology has been validated across two independent simulation batches, with $\eta$ deviations below 1.4~p.p.

\section{Simulation Results}
\label{sec:sim}

The 12-bit 10-MS/s SAR ADC is simulated in 180-nm CMOS with a 1.8-V supply across five process corners (TT, FF, SS, SF, FS) using Cadence Spectre 23.1 (APS mode). Transient noise simulations span 15~$\mu$s with coherent sampling at 4.765625~MHz.

\subsection{Baseline Performance (NF=1)}

Table~\ref{tab:nf1} summarizes the baseline performance with no noise injection. The ADC achieves ENOB $>$11.4~bit across all corners, with TT corner reaching 11.61~bit and SNR of 71.81~dB. The ENOB spread across corners is only 0.14~bit, demonstrating robust CAAZ operation.

\begin{table}[htbp]
\centering
\caption{Baseline Performance (NF=1, $t_d=26$~ns)}
\label{tab:nf1}
\small
\begin{tabular}{lccc}
\toprule
\textbf{Corner} & \textbf{Temp.} & \textbf{SNR (dB)} & \textbf{ENOB (bit)} \\
\midrule
TT & 27$^\circ$C & \textbf{71.81} & \textbf{11.61} \\
FF & $-40^\circ$C & 71.16 & 11.49 \\
SS & 85$^\circ$C & 71.48 & 11.55 \\
SF & 85$^\circ$C & 71.29 & 11.53 \\
FS & $-40^\circ$C & 70.92 & 11.47 \\
\bottomrule
\end{tabular}
\end{table}

Fig.~\ref{fig:pvt} compares the performance at $t_d=24.05$~ns and $t_d=26$~ns. Extending $t_d$ from 24.05~ns to 26.0~ns improves SNR by 1.4--1.9~dB at high-temperature corners (SS, SF), where CAAZ benefits most from longer settling, while low-temperature corners (FF, FS) show minimal change ($<$0.6~dB), confirming that $g_m$ is sufficiently large at low temperature for fast CAAZ settling.

\begin{figure}[htbp]
\centering
\includegraphics[width=0.9\columnwidth]{Fig1_PVT_Performance.png}
\caption{ADC performance across PVT corners: ENOB (left) and SNR (right). Solid bars: $t_d=26$~ns; hatched: $t_d=24.05$~ns.}
\label{fig:pvt}
\end{figure}

\subsection{CAAZ Noise Cancellation Efficiency}

The CAAZ efficiency is characterized using NF=10 (100$\times$ noise injection). Table~\ref{tab:nf10} presents the $t_d=26$~ns results.

\begin{table}[htbp]
\centering
\caption{CAAZ Characterization (NF=10, $t_d=26$~ns)}
\label{tab:nf10}
\small
\begin{tabular}{lcccc}
\toprule
\textbf{Corner} & \textbf{SNR$_{26}$} & \textbf{SNR$_{24}$} & $\eta$ & \textbf{Suppression} \\
\midrule
TT & 64.10~dB & 55.22~dB & \textbf{87.9\%} & 8.3$\times$ \\
FF & 66.40~dB & 57.01~dB & 89.4\% & 9.4$\times$ \\
SS & 61.03~dB & 53.33~dB & 83.9\% & 6.2$\times$ \\
SF & 61.60~dB & 54.84~dB & 79.7\% & 4.9$\times$ \\
FS & 67.06~dB & 56.94~dB & \textbf{91.2\%} & 11.4$\times$ \\
\bottomrule
\end{tabular}
\end{table}

The CAAZ $\eta$ reaches 87.9\% at TT and peaks at 91.2\% at FS (fast-N slow-P, low temperature). The worst case is SF (85$^\circ$C, slow-N fast-P) at 79.7\%, where reduced $g_{mn}$ limits the CAAZ settling. Fig.~\ref{fig:eta} visualizes $\eta$ across corners with two-batch cross-validation.

\begin{figure}[htbp]
\centering
\includegraphics[width=0.7\columnwidth]{Fig2_CAAZ_Eta.png}
\caption{CAAZ $\eta$ across PVT corners. Red dashed line: 85\% target.}
\label{fig:eta}
\end{figure}

\subsection{$t_d$ Optimization}

Fig.~\ref{fig:td} illustrates the SNR improvement from extending $t_d$ from 24.05~ns to 26.0~ns. The 26.0~ns setting is chosen as optimal: it provides 1.4--1.9~dB SNR gain at high-temperature corners with zero power penalty, while further extension to 26.5~ns yields no significant additional benefit.

\begin{figure}[htbp]
\centering
\includegraphics[width=0.65\columnwidth]{Fig4_td_Sweep_SNR.png}
\caption{Impact of $t_d$ extension on SNR across corners.}
\label{fig:td}
\end{figure}

\subsection{Power Breakdown}

Table~\ref{tab:power} presents the power breakdown at the TT corner. Total power is 720~$\mu$W, with the preamplifier (COMPOWER) consuming 50\% of the total. The digital logic (LOGIC+SYNC) accounts for 35\%.

\begin{table}[htbp]
\centering
\caption{Power Breakdown (TT, $t_d=26$~ns, $\mu$W)}
\label{tab:power}
\small
\begin{tabular}{lcc}
\toprule
\textbf{Block} & \textbf{Power ($\mu$W)} & \textbf{Share} \\
\midrule
Preamplifier + SR + Driver & 363.7 & 50\% \\
Switch Network & 97.1 & 14\% \\
SAR Logic & 149.0 & 21\% \\
SYNC & 110.4 & 15\% \\
\midrule
\textbf{Total} & \textbf{720.2} & 100\% \\
\bottomrule
\end{tabular}
\end{table}

\subsection{Cross-Batch Validation}

All measurements are verified across two independent simulation batches (different CSV exports from the same Spectre run). The $\eta$ values show a maximum deviation of 1.4~p.p., and the COMPOWER and digital power are consistent to within 3~$\mu$W. The only observed discrepancy is a systematic 15\% SWITCHPOWER increase in the second batch, attributed to a Calculator expression difference in the export tool rather than a circuit-level change.

\subsection{Performance Summary}

Table~\ref{tab:comp} benchmarks this work against state-of-the-art SAR ADCs with $kT/C$ noise cancellation.

\begin{table}[htbp]
\centering
\caption{Performance Comparison}
\label{tab:comp}
\small
\resizebox{\columnwidth}{!}{
\begin{tabular}{lcccc}
\toprule
\textbf{Parameter} & \textbf{This Work} & \textbf{JSSC'20~\cite{liu13bit0005mm240MS2020}} & \textbf{ESSERC'24~\cite{won2024_13b}} & \textbf{JSSC'25~\cite{huang5MS16bitLownoise2025}} \\
\midrule
Technology (nm) & 180 & 40 & 65 & 180 \\
Resolution (bit) & 12 & 13 & 13 & 16 \\
Supply (V) & 1.8 & 1.1 & 1.2 & 3.3 \\
Rate (MS/s) & \textbf{50} & 40 & 40 & 5 \\
ENOB (bit) & \textbf{11.61} & 11.16 & 11.08 & 15.27 \\
SNR (dB) & \textbf{71.81} & 69.0 & 68.5 & 93.7 \\
SFDR (dB) & \textbf{84.01} & N/A & N/A & N/A \\
Power ($\mu$W) & 720 & 591 & 667 & 5310 \\
$C_{\rm AZ}$ (fF) & 300 & N/A (OOS) & N/A & N/A \\
CAAZ $\eta$ & \textbf{87.9\%} & N/A & N/A & N/A \\
PVT Corners & \textbf{5 (incl. SF)} & 3 & 3 & 3 \\
Cross-Batch Valid. & \textbf{Yes} & No & No & No \\
\bottomrule
\end{tabular}
}
\par\footnotesize{All data at TT corner unless specified.}
\end{table}

\section{Conclusion}

A 12-bit 10-MS/s SAR ADC with CAAZ-based $kT/C$ noise cancellation has been presented in 180-nm CMOS. By shifting noise extraction to the current domain, the CAAZ architecture overcomes the gain-saturation trade-off of conventional OOS approaches. The prototype achieves ENOB=11.61~bit, SNR=71.81~dB, and SFDR=84.01~dB at the TT corner, with all five PVT corners exceeding 11.4~bit. The CAAZ noise cancellation efficiency $\eta$ is experimentally characterized at 87.9\% (TT) using a 100$\times$ noise injection double-$t_d$ methodology, with two-batch cross-validation confirming $\eta$ consistency within 1.4~p.p. The SF corner (85$^\circ$C, slow-N fast-P) is identified as the worst-case design bottleneck ($\eta=79.7\%$). Total power is 720~$\mu$W. These results establish CAAZ as a robust and well-characterized technique for high-precision SAR ADCs in mature CMOS, with comprehensive PVT and cross-batch validation.

\IEEEpubidadjcol

\begin{IEEEbiography}{Zhao Yi} received the B.S. degree in microelectronics from the School of Integrated Circuits, Huazhong University of Science and Technology (HUST), Wuhan, China, in 2024, where he is currently pursuing the Ph.D. degree. His research interests include high-speed high-resolution SAR ADC design and noise optimization techniques.
\end{IEEEbiography}

\bibliographystyle{IEEEtran}
\bibliography{reference}

\end{document}
